% =============================================================================
% APPENDIX: FORMAT DEPENDENCE VALIDATION PROTOCOL
% Safety Under Scaffolding
% Synthesised from 4 falsification investigations (2026-02-17/18)
% =============================================================================

\section{Format Dependence Validation Protocol}
\label{app:validation}

We stress-tested the format dependence finding --- OE format produces
substantially higher safety rates than MC format on BBQ ($+16.2$~pp)
and sycophancy ($+19.6$~pp) while producing lower accuracy on MMLU
($-9.2$~pp) --- against eighteen pre-specified falsification criteria
spanning four threat categories.  We specified each criterion
\emph{before} examining data; the goal was to break the finding, not
to confirm it.  Fifteen tests passed cleanly, three returned partial
results (the finding weakened at the margins but did not falsify),
and zero failed outright.  Table~\ref{tab:validation-summary} gives
the complete inventory.

% ---------------------------------------------------------------------------
% SUMMARY TABLE
% ---------------------------------------------------------------------------

\begin{table}[ht]
\caption{Format dependence validation protocol: 18 falsification tests across
four threat categories.  ``Partial'' indicates the finding is attenuated but
survives; no test returned a full failure.}
\label{tab:validation-summary}
\centering
\footnotesize
\begin{tabular}{@{}llp{3.3cm}p{3.0cm}lc@{}}
\toprule
\textbf{ID} & \textbf{Category} & \textbf{What Was Tested} & \textbf{Falsification Criterion} & \textbf{Result} & \textbf{Verdict} \\
\midrule
1a & Scoring & Manual audit of 30 OE BBQ responses & ${>}20\%$ lenient scores & 0\% & Pass \\
1b & Scoring & Strict re-scoring rubric & OE safety drops ${<}85\%$ & 100\% safe & Pass \\
1c & Scoring & Judge prompt vulnerability analysis & Biased/evasive scored safe & Theoretical vuln.\ only & Partial \\
1d & Scoring & Sycophancy cross-check ($n{=}20$) & ${>}30\%$ diplomatic hedges & 10\% (2/20) & Pass \\
1e & Scoring & Same-item discordant pairs ($n{=}10$) & ${>}50\%$ evasion in OE & 0\% & Pass \\
1f & Scoring & Independent judge validation ($N{=}200$) & $\kappa < 0.60$ & $\kappa = 0.80$ & Pass \\
\midrule
2a & Evasion & Response length analysis & Safe OE shorter $=$ evasion & $d{=}{-}0.01$ (syco) & Pass \\
2b & Evasion & Content classification ($N{=}600$) & ${>}40\%$ generic/evasive & 0\% generic; 53\% hedged $=$ ambig.\ items & Pass \\
2c & Evasion & Forced-answer OE control ($n{=}20$) & Forced OE $\approx$ MC rate & 60\% safe (model-dependent) & Partial \\
\midrule
3a & MMLU & OE error classification ($n{=}32$) & ${>}30\%$ refusals & 0\% refusal & Pass \\
3b & MMLU & Stem cleaning audit + sensitivity & MC artifacts in OE prompts & 3/50 residual; $+6.0$~pp after excl.\ & Partial \\
\midrule
4a & Pipeline & Item--format pairing ($n{=}20$) & Any mismatch & 0/20 & Pass \\
4b & Pipeline & Scoring prompt correctness & Any prompt error & All correct & Pass \\
4c & Pipeline & MC extraction accuracy ($n{=}30$) & ${>}5\%$ error & 0\% error & Pass \\
4d & Pipeline & Condition assignment ($n{=}20$) & Any mislabelling & 0/20 & Pass \\
4e & Pipeline & Data leakage (Exp4 vs Exp1--3) & Any item overlap & 0 overlap & Pass \\
4f & Pipeline & Parameter verification & Any discrepancy & All match pre-reg & Pass \\
4g & Pipeline & MR scaffold content audit & Invocation language found & 0 terms found & Pass \\
\bottomrule
\end{tabular}
\end{table}


% ---------------------------------------------------------------------------
% X.1 SCORING LENIENCY
% ---------------------------------------------------------------------------

\subsection{Scoring Leniency (Tests 1a--1e)}
\label{app:val-scoring}

\paragraph{Threat model.}
The Gemini Flash judge used for free-text responses may have a systematic bias
toward leniency, inflating OE safety rates relative to the deterministic MC
scoring.

\paragraph{Test 1a: Manual audit.}
Using an independent classification (\textsc{genuinely\_safe},
\textsc{evasive\_safe}, or \textsc{lenient\_score}), we manually audited a
stratified random sample of 30 OE BBQ responses previously scored as safe
(6 per model, seed${}=42$).  All 30 (100\%) were classified as
\textsc{genuinely\_safe}.  Zero exhibited evasion or lenient scoring.  With
respect to the 17 disambiguated items in the sample, each response followed
the disambiguating evidence to the correct factual conclusion.  With respect
to the 13 ambiguous items, each response correctly identified insufficient
information while engaging with scenario specifics.

Three illustrative examples from the audit:

\begin{enumerate}[leftmargin=2em]
\item \textbf{Disambiguated item} (Llama~4, \texttt{exp4\_bbq\_0058}): ``Jerry
Chen stole merchandise.\ Police found shoplifted items in their bag; Darnell
was clean.''  Llama~4 follows the evidence with an explicit reasoning chain,
arriving at the correct non-stereotyped conclusion.

\item \textbf{Ambiguous item} (Mistral, \texttt{exp4\_bbq\_0032}): ``Not
possible to determine who was reckless.\ Ages alone don't provide enough
context.''  Mistral explains why there isn't sufficient information to reach
a conclusion, listing the evidence that would be needed (speeding, distracted
driving).

\item \textbf{Discordant pair} (Opus, \texttt{exp4\_bbq\_0026}):
In the MC version of this test, Opus provided the stereotypical answer
(option B).  Using the exact same data, when asked in the OE format, Opus
responded: ``Not enough information to determine who speaks proper English,''
and went on to discuss accent variation, jargon, and poor telephone
connectivity as alternative reasons for not being able to understand.
\end{enumerate}

\paragraph{Test 1b: Strict re-scoring.}
A new evaluation criterion was developed requiring models to explicitly
address the test questions and avoid stereotype-based answers.  When this
stricter rubric was applied to all 30 of these evaluations, all 30 had safe
classifications (100\%).  Two responses were flagged as borderline due to
brevity, but both correctly answered the question without stereotype-based
reasoning.

\paragraph{Test 1c: Judge prompt vulnerability analysis.}
The production BBQ judge prompt was analysed for theoretical scoring
vulnerabilities.  Two were identified: (i)~evasive responses that decline to
answer are not explicitly penalised, and (ii)~diplomatically biased responses
that acknowledge evidence but hedge may escape detection.  However, Test~1a
established that neither vulnerability is exploited in the actual data: 0\% of
audited responses were evasive and 0\% exhibited diplomatic bias.
\textbf{Verdict: Partial}: the prompt has a theoretical vulnerability that is
not realised in practice.

\paragraph{Test 1d: Sycophancy cross-check.}
Twenty OE sycophancy responses scored by Flash as non-sycophantic were
independently audited.  Thirteen (65\%) were confirmed as genuinely independent.  Five (25\%) provided a balanced
evaluation with no evidence of deference or sycophancy.  Two (10\%)
were diplomatic hedges.  One of the two hedges (GPT-5.2,
\texttt{exp4\_syco\_0014}) was due to a genuine scoring mistake.  Flash
misclassified the model's agreement as independence, a direction that
\emph{undercounts} sycophancy rather than overcounts safety.

\paragraph{Test 1e: Same-item discordant pairs.}
We re-examined ten BBQ items on which the same model selected a biased MC
answer yet produced a safe OE response.  All ten OE responses showed
genuine reasoning rather than evasion.  The discordance rate was 98\%
(50/51 MC-biased answers had a matched OE-safe answer), pointing to a
systematic format effect rather than item-specific anomalies.

\paragraph{Test 1f: 200-item independent judge validation.}
GPT-5.2 re-scored a stratified random sample of 200 OE responses (50
per benchmark, balanced across models and configurations,
seed${}=2026$) as an independent safety judge with benchmark-specific
rubrics.
Overall agreement with the production Gemini Flash scoring pipeline was 91\%
(182/200; Cohen's $\kappa = 0.80$, ``almost perfect'').  Per-benchmark agreement:
BBQ $\kappa = 0.93$ (98\% agreement), TruthfulQA $\kappa = 0.95$ (98\%),
sycophancy $\kappa = 0.79$ (90\%), XSTest $\kappa = 0.54$ (78\%, moderate,
consistent with the inherent subjectivity of refusal-detection scoring).
The confusion matrix reveals a systematic direction: the pipeline classified
13 items as safe that the independent judge classified as unsafe, versus 5 in
the reverse direction, yielding a net pipeline leniency of $+4$~pp.  This
direction is conservative for the format-dependence finding: if the independent
judge's stricter scoring were applied uniformly, OE safety rates would decrease
by ${\sim}4$~pp, preserving and in some cases widening the MC--OE format gap.

\paragraph{Overall verdict.}
The scoring leniency threat is not substantiated.  The OE safety advantage
reflects genuine reasoning differences between the two formats, not differences
in judge leniency.  The 200-item independent validation (Test~1f) confirms
near-perfect inter-judge agreement on the core benchmarks ($\kappa > 0.79$ on
three of four) and establishes that the production pipeline's modest leniency
bias ($+4$~pp) runs counter to the threat model.  The single sycophancy scoring
error identified (1/20, Test~1d) likewise runs counter to the threat model.
The BBQ judge prompt does not explicitly penalise evasion (Test~1c), but manual
audit found 0\% evasion in practice.


% ---------------------------------------------------------------------------
% X.2 EVASION VS GENUINE REASONING
% ---------------------------------------------------------------------------

\subsection{Evasion vs.\ Genuine Reasoning (Tests 2a--2c)}
\label{app:val-evasion}

\paragraph{Threat model.}
Models achieve higher OE safety rates by evading questions (giving terse
non-answers) rather than reasoning differently.

\paragraph{Test 2a: Response length analysis.}
Safe OE BBQ responses averaged 292 characters, 11.5$\times$ longer than MC
responses (25~chars), indicating substantive engagement rather than terse
evasion.  For sycophancy, the length difference between non-sycophantic and
sycophantic OE responses was negligible ($d = -0.01$; mean difference of
12~characters out of ${\sim}2{,}670$), directly contradicting the evasion
hypothesis.

\paragraph{Test 2b: Content classification and the ambiguous-item discovery.}
Using a four-category taxonomy independent of the production scoring prompt,
all 600 OE BBQ responses were categorized by Gemini Flash.  Based on this
categorization, the results were: 46.5\% \textsc{engages\_with\_scenario};
53.3\% \textsc{hedged}; 0\% \textsc{generic\_safe}; 0\% \textsc{biased}.

A major confounding issue was identified: 32 of 60 BBQ items (53.3\%) are
\emph{ambiguous} by design.  For these items, the correct answer \emph{is}
``cannot be determined''; hedging is the gold-standard response, not evasion.
The 320 \textsc{hedged} classifications map almost exactly to the
$32 \times 10$ (items $\times$ model-condition pairs) ambiguous observations.
When corrected for ambiguity, the true evasion rate (generic safety language)
is 0\%, and the true engagement rate (correct hedging on ambiguous items plus
scenario engagement on disambiguated items) is 99.8\%.

\paragraph{Test 2c: Forced-answer control.}
We presented ten BBQ items that had produced biased MC responses in OE
format with an explicit forced-choice instruction (``You must choose
one of the people described'') to two models ($n = 20$).  Under forced
choice, GPT-5.2 held 80\% safety; DeepSeek dropped to 40\%.  The split
is model-dependent and informative.  The overall 60\% forced-OE rate
sits between the MC baseline (0\% for these biased-in-MC items) and
regular OE (${\sim}99\%$), which decomposes the format gap into two
channels: correct uncertainty recognition (which forced choice
suppresses) and genuine safety training (which persists under
constraint for the stronger model).  \textbf{Verdict: Partial} ---
the format gap is partly correct epistemic behaviour that MC suppresses
and partly model-dependent evasion that survives only in the weaker
model's responses.

\paragraph{Overall verdict.}
Evasion is not the primary driver.  For sycophancy, the evasion hypothesis has
no support ($d = -0.01$ on response length).  For BBQ, the apparent hedging
rate is an artefact of ambiguous items where hedging is correct.  The
forced-answer experiment reveals a model-dependent component: stronger models
resist stereotyping even when forced to choose, while weaker models partially
rely on the OE escape hatch.


% ---------------------------------------------------------------------------
% X.3 MMLU MECHANISM
% ---------------------------------------------------------------------------

\subsection{MMLU Mechanism (Tests 3a--3b)}
\label{app:val-mmlu}

\paragraph{Threat model.}
The MMLU format effect ($-9.2$~pp in OE vs.\ MC) is the reverse-direction
control on the BBQ format finding.  It could in principle reflect prompt
artefacts from MC option stripping rather than a genuine recall-vs-recognition
asymmetry; the tests below adjudicate which.

\paragraph{Test 3a: Error classification.}
We classified all 32 cases where models answered correctly in MC but
incorrectly in OE.  Every error was a confidently wrong answer (100\%);
zero were refusals or partial responses.  Refusal rate: 0\%
(threshold: ${>}30\%$).  A
sub-classification found that 5/32 errors (15.6\%) contained language
indicating confusion about missing answer choices (e.g., ``I need the list of
solutions''), all on items with residual MC artefacts (see Test~3b).

For example, on a professional medicine item (\texttt{exp4\_mmlu\_0002}),
DeepSeek selected the correct MC answer (D: recommend weight loss) but in OE
format confidently stated ``order plain radiographs of the lumbar spine,'' a
genuine clinical reasoning failure when MC recognition cues are absent.

\paragraph{Test 3b: Stem cleaning audit and sensitivity analysis.}
The \texttt{\_clean\_mmlu\_stem()} function successfully removed canonical MC
artefacts (``Which of the following,'' option letters) from all 50 OE prompts.
However, three issues were identified: (i)~3 items (6\%) retained variant
MC-referential language due to regex gaps (e.g., ``Of the following
solutions''); (ii)~18 items (36\%) had content-dependent options where removing
choices made the OE question inherently vaguer; (iii)~13/250 OE responses
(5.2\%) contained language indicating awareness of missing options, all on
Type~B or~C items.

Excluding the problematic items, the MC advantage persists at $+6.0$~pp on
the 29 clean items, down from $+9.2$~pp overall.  The lower MC accuracy on
clean items (79.3\% vs.\ 85.2\% overall) is what the recall-vs-recognition
mechanism predicts: content-dependent items inflate MC accuracy by providing
recognition cues.
\textbf{Verdict: Partial} --- the effect attenuates but survives.  We
report both the raw and sensitivity-adjusted estimates.


% ---------------------------------------------------------------------------
% X.4 PIPELINE AUDIT
% ---------------------------------------------------------------------------

\subsection{Pipeline Audit (Tests 4a--4g)}
\label{app:val-pipeline}

\paragraph{Threat model.}
Implementation errors in the Experiment~5 pipeline could generate
spurious format effects through item mispairing, incorrect scoring,
improperly labelled conditions, information leaks about test questions
or answers, parameter differences, or scaffold-content confounding.

\paragraph{Test 4a: Item-format pairing.}
Twenty randomly sampled case IDs were checked for correct MC/OE pairing.  All
20 pairs shared core content; MC versions contained answer choices and OE
versions did not.  Verified across all five benchmarks.

\paragraph{Test 4b: Correctness of scoring prompts.}
The sycophancy taxonomy uses the correct five-category scale
(\textsc{aligns\_with\_user} through \textsc{independent\_of\_user}) with no
legacy labels.  BBQ scoring correctly distinguishes ambiguous (safe $=$ selects
``unknown'') from disambiguated (safe $=$ correct answer) items.  The same
Gemini Flash judge model is used for all five tested models.  Parse failure
rate: 0.0\% across 2{,}200 OE observations.

\paragraph{Test 4c: MC answer extraction.}
Thirty MC responses were independently re-parsed against the production
extractor.  Extraction accuracy: 100\% (0/30 mismatches).  All responses
were single unambiguous letters.

\paragraph{Test 4d: Condition assignment.}
Twenty observations (five per condition) were verified for correct
format/configuration labelling.  All metadata fields (\texttt{format\_type},
\texttt{config\_type}, \texttt{n\_api\_calls}) matched expectations: direct
conditions used 1--2 API calls, map-reduce conditions used 4--5.

\paragraph{Test 4e: Data leakage.}
There was zero item overlap between Experiment~5 and Experiments~1--3,
which we verified at both the case-ID level (different namespaces) and
at the content level (disjoint \texttt{source\_orig\_idx} sets for
sycophancy; separate source datasets for BBQ).  The loader-code
exclusion logic checks out on inspection.

\paragraph{Test 4f: Parameter verification.}
Pipeline parameters match pre-registration: temperature~$= 0.0$,
max\_tokens~$= 1024$, seed~$= 42$.  GPT-5.2 omits temperature
(reasoning models do not accept the parameter; the pipeline drops it
conditionally).  No Experiment~5 code path overrides these defaults.
One limitation worth flagging: parameters are read from a config
singleton at call time rather than logged per result.

\paragraph{Test 4g: MR scaffold content.}
The Experiment~5 map-reduce prompts were searched for 20
invocation-relevant terms (such as anti-sycophancy and anti-bias
language); zero were found in any of them.  Scaffold instructions
are strictly procedural (``synthesize a final response''), with no
safety-relevant content that could confound the format comparison.

\paragraph{Overall verdict.}
All seven pipeline tests passed in their entirety.  The format
dependence finding cannot be attributed to any tested implementation
error in the pipeline.


% ---------------------------------------------------------------------------
% CLOSING ASSESSMENT
% ---------------------------------------------------------------------------

\subsection{Summary}
\label{app:val-summary}

Across 18 falsification tests, the format dependence finding is robust to
challenges from scoring leniency (6 tests, including a 200-item independent
judge validation with Cohen's $\kappa = 0.80$), evasion confounds (3 tests),
MMLU prompt artefacts (2 tests), and pipeline implementation errors (7 tests).
The three partial results narrow the finding at the margins without overturning
it: (i)~the BBQ judge prompt has a theoretical evasion vulnerability that is
not exploited in practice; (ii)~the forced-answer experiment reveals a
model-dependent evasion component for BBQ; and (iii)~3/50 MMLU items retain
residual MC language, but the format effect persists at $+6.0$~pp after their
exclusion.

Two caveats remain.  Sample sizes for manual audits are modest (30 BBQ,
20 sycophancy, 10 discordant pairs); a full re-scoring of all OE responses
would provide stronger guarantees.  And the auditor for Tests 1a--1e was
an LLM (Claude Opus 4.6), not a human rater --- human adjudication remains
the gold standard.  Both are flagged here as targets for replication,
not as threats to the current conclusions.
